% Generated draft: requires mathematical and source review.
\documentclass[11pt]{article}
\usepackage[margin=25mm]{geometry}
\usepackage{fontspec}
\setmainfont{Latin Modern Roman}
\usepackage{amsmath,amssymb,mathtools,unicode-math}
\setmathfont{Latin Modern Math}
\usepackage{enumitem,needspace,xurl,hyperref}
\setlength{\parindent}{0pt}
\setlength{\parskip}{5pt}
\setlength{\emergencystretch}{3em}
\newcommand{\N}{\mathbb{N}}
\newcommand{\Z}{\mathbb{Z}}
\newcommand{\Q}{\mathbb{Q}}
\newcommand{\R}{\mathbb{R}}
\newcommand{\C}{\mathbb{C}}
\newcommand{\F}{\mathbb{F}}
\newcommand{\A}{\mathbb{A}}
\newcommand{\PP}{\mathbb{P}}
\newcommand{\E}{\mathbb{E}}
\newcommand{\eps}{\varepsilon}
\newcommand{\dd}{\,\mathrm{d}}
\begin{document}
\section*{UnsolvedMath ID: 20002057}
\textbf{Source problem code:} catalogue entry
\subsection{Definitions and mathematical statement}
Problem 33 asks for general or inductive formulae for the operators whose OCR name is ``§2`''. The aid identifies them as the conformal tractor Laplacians \(\Box_{2k}\) and attributes a general effective curvature-correction algorithm to Gover--Peterson. The record does not include that full algorithm or its citation. Its concrete result is an Einstein-scale factorization with the following explicit range and conventions.

Let \(n=2p\geq4\), choose an Einstein representative \(g\) of a Riemannian conformal structure, and set
\[
 J=\frac{R_g}{2(n-1)},\qquad
 \operatorname{Ric}_g=\frac{2(n-1)}nJg.
\]
Then \(J\) is constant. On the tractor bundle under consideration let \(\Delta^{\mathcal T}=\nabla^a\nabla_a\) be the coupled rough Laplacian and put \(A=-\Delta^{\mathcal T}\). An Einstein scale trivializes density weights, allowing the factors to be viewed as operators on one underlying bundle. The aid reports, for the canonical weighted tractor operator,
\[
 \Box_{2k}=\prod_{m=1}^k
 \left(A+\frac{2J}{n}(p-m)(p+m-1)\right),
 \qquad1\leq k\leq p-1.
\]
With weights restored, the usual assignment is
\(\mathcal V[k-p]\to\mathcal V[-k-p]\) for the relevant tractor bundle \(\mathcal V\). The factors commute in the chosen scale because they are polynomials in the same \(A\) with constant coefficients. This includes Einstein metrics with nonzero Weyl curvature, not just conformally flat examples.

The corresponding fixed-dimension one-step algorithm is
\[
 P_0=I,\qquad
 P_k=\left(A+\frac{2J}{n}(p-k)(p+k-1)\right)P_{k-1},
 \quad1\leq k\leq p-1.
\]
It is a polynomial recursion after scale identification; it is not an unqualified composition between differently weighted intrinsic operators. No extension to \(k\geq p\) is inferred from the phrase ``all orders.'' For a conformally Einstein structure one must work where its Einstein representative is defined and nondegenerate; a vanishing almost-Einstein scale requires separate treatment.

For the standard tractor bundle at \(k=p-1\), the operator has weights
\(\mathcal T[-1]\to\mathcal T[1-n]\). The aid's coefficient checks are
\[
 [A^{p-2}]\Box_{n-2}=\frac{(n-1)(n-2)}6J,
 \qquad
 [A^0]\Box_{n-2}=(n-2)!\left(\frac{2J}{n}\right)^{p-1}.
\]
Indeed the shifts sum to the first coefficient and their product gives the second. For \(n=4\), both refer to the same constant coefficient and \(\Box_2=A+J\); for \(n=6\),
\[
 \Box_4=A^2+\frac{10}{3}JA+\frac83J^2.
\]
When \(J=0\), the factorization reduces to \(A^k\). These are tests of an Einstein restriction, not a complete verification of a proposed formula on arbitrary non-Einstein conformal geometries.

The extraction then begins the next problem's setup: a natural tensor-density \(N_g\) may have the finite affine conformal law
\[
 N_{e^{2\omega}g}=N_g+L_g\omega,
\]
with a universal linear differential operator in \(\omega\), recorded as equation (43). Its target weights and full next question are absent. This background does not replace the operator-formula task or supply its missing general curvature corrections.
\subsection{Short English statement}
The aid reports a general tractor-calculus algorithm and gives a closed factorization in an Einstein scale for the stated subcritical orders. A fixed-dimension recursion and two critical-operator coefficient checks follow. The formula is an Einstein restriction, with weights and order range preserved, rather than a full arbitrary-geometry expansion.
\subsection{Sources}
\begin{itemize}
\item Permanent UnsolvedMath ID \(20002057\), exact public code catalogue entry; local registry \texttt{aim-geometry-notes.json}, \texttt{notes[394]}, canonical tag problem.
\item AIM workshop ``Conformal structure in geometry, analysis, and physics,'' Problem 33 and adjoining equation (43): \url{https://aimath.org/WWN/confstruct/confstruct.pdf}. No PDF was accessed. The supplied verification note identifies OCR damage and the opening paragraph of the next problem.
\item The aid attributes the effective general curvature-correction construction to Gover--Peterson, without a full title, theorem reference or algorithm transcription. No original literature entry, proposer or year is supplied.
\item Research attempt 1, ``General tractor Laplacians and their Einstein factorization,'' labelled partial\_result, partially\_solved, proved and low novelty confidence. The factorization and coefficient identities are supplied research reports rather than independently checked source results.
\end{itemize}
\end{document}
